% Product Management – Midterm Study Guide — EU business casual
% Compile with: pdflatex study-guide.tex

\documentclass[8pt,landscape,a4paper]{article}
\usepackage[utf8]{inputenc}
\usepackage[T1]{fontenc}
\usepackage[left=0.8cm,right=2.6cm,top=0.8cm,bottom=0.8cm,marginparwidth=2cm,marginparsep=0.3em]{geometry}
\usepackage{multicol}
\usepackage{marginnote}
\usepackage{enumitem}
\setlist{nosep,leftmargin=*}
\usepackage{booktabs}
\usepackage{parskip}
\usepackage{microtype}
\usepackage{xcolor}
\usepackage{fancyhdr}
\definecolor{eublue}{RGB}{0,51,153}
\definecolor{charcoal}{RGB}{51,51,51}
\definecolor{softgrey}{RGB}{240,240,240}
\definecolor{notegrey}{RGB}{90,90,90}
\newcommand{\handnote}[1]{\marginnote{\raggedright\footnotesize\color{notegrey}\textit{#1}\vspace{0.4em}}}
\newcommand{\hlpink}[1]{\textcolor{charcoal}{\textbf{#1}}}
\newcommand{\hlpeach}[1]{\textcolor{eublue}{\textbf{#1}}}
\setlength{\parskip}{0.2em}
\setlength{\columnsep}{0.5em}
\pagestyle{fancy}
\fancyhf{}
\fancyfoot[L]{\raisebox{0.8em}{\small\color{notegrey} Product Mgmt}}
\fancyfoot[R]{\raisebox{0.8em}{\small\color{notegrey}\thepage}}
\fancypagestyle{plain}{\fancyhf{}\fancyfoot[L]{\raisebox{0.8em}{\small\color{notegrey} Product Mgmt}}\fancyfoot[R]{\raisebox{0.8em}{\small\color{notegrey}\thepage}}}
\raggedright

\begin{document}
\begin{center}
\vspace{0.2em}
{\normalsize\bfseries\color{eublue} Product Management -- Midterm Study Guide}\\[0.2em]
\footnotesize\textcolor{notegrey}{Use for quick recall; in exam: define terms, give one example, tie to value creation/capture.}\\[0.35em]
\noindent\colorbox{softgrey}{\parbox{0.92\textwidth}{\centering\footnotesize\color{charcoal}
\textbf{Key:} \textbf{Definitions} \quad \textcolor{eublue}{\textbf{Frameworks \& tips}} \quad \textit{Margin notes}
\vspace{0.25em}}}
\end{center}
\vspace{-0.4em}

\begin{multicols}{2}

\section*{\color{eublue}Unit 0 -- Course \& Big Ideas}\handnote{Creation $\neq$ capture.}
\begin{itemize}[leftmargin=*,nosep]
\item \textbf{Product management} = Creating AND capturing value. Great product is not enough.
\item \textbf{Goal}: CSDA (customer-focused, strategic, data-driven, action-oriented).
\item Key insight: ``Having a great product isn't enough''---you need market access to capture value.
\item Valuation: asset-based (what you have) vs.\ market-based (market size, potential customers, access).
\item \textbf{Don't build and hope they come}---understand before you build.\handnote{Community before product.}
\end{itemize}

\section*{\color{eublue}Unit 1 -- Value Creation vs.\ Capture}\handnote{Need both. Teece.}
\begin{itemize}[leftmargin=*,nosep]
\item \textbf{\hlpink{Value creation}} = Making something useful. Total value for users.
\item \textbf{\hlpink{Value capture}} = How much money you make. You need both.
\item \textbf{\hlpink{Market access}} = Distribution, channels. Teece: market access $\to$ value capture.
\item \textbf{Porter value chain}: Primary (production, delivery) + support (HR, tech). Value = resources + activities $\to$ users.
\item \textbf{Takeda case}: Spanish research 15K EUR $\to$ Takeda 250M EUR/year. Better market access, distribution, scale.\handnote{Classic example!}
\item \textbf{Value capture mechanisms}: Patents, branding, market access, co-specialized assets.
\item \textbf{Co-specialized assets}: Specialized to your business; provide market access. Branding, secrecy, distribution, customer relationships.
\item Patents: protect innovation; help capture even after selling. Avoid over-patenting.
\item European challenge: smaller market reach limits value capture vs.\ global players.
\end{itemize}

\subsection*{Value creation vs.\ capture}
\begin{center}
\footnotesize
\begin{tabular}{@{}lll@{}}
\toprule
\textbf{Aspect} & \textbf{Creation} & \textbf{Capture} \\
\midrule
What & Total value for users & Money you make \\
Depends on & Product, usefulness & Market access, patents, brand \\
\bottomrule
\end{tabular}
\end{center}

\section*{\color{eublue}Unit 2 -- Community \& Needs Assessment}\handnote{Before you build.}
\begin{itemize}[leftmargin=*,nosep]
\item \textbf{Community before product}: Understand what they need first. Don't build and hope they come.
\item \textbf{Needs assessment}: Find connections. Assess what community actually needs.
\item \textbf{Willingness to pay}: Test early. If no one pays, product isn't solving real problem.
\item \textbf{Testing as investing}: Testing is part of development, not separate.
\item Social app failure: spent 500K, no one wanted to pay. Lesson: build community first.
\item Google model: free service, value capture through advertising.
\end{itemize}

\section*{\color{eublue}Unit 3 -- Cost \& Production}\handnote{Cost $\neq$ price.}
\begin{itemize}[leftmargin=*,nosep]
\item \textbf{Cost $\neq$ price}. Low cost doesn't mean low price. You choose price.
\item \textbf{Fixed costs}: Building, insurance. Don't change with production.
\item \textbf{Variable costs}: Materials, wages. Change with production.
\item \textbf{Economies of scale}: Larger production = lower cost per unit. Spread fixed costs.
\item \textbf{Standardization}: Cheaper, faster. Mass production requires it.
\item \textbf{Specialization}: Workers doing one task = more efficient. Learning by doing.
\item \textbf{Henry Ford Model T}: 12h $\to$ 2.5h per car; \$850 $\to$ \$260. Assembly line, standardized parts.
\item \textbf{Mass production requires scale}: Need enough volume to cover fixed costs.
\end{itemize}

\section*{\color{eublue}Unit 4 -- BMC \& Product Design}\handnote{Validate before you build.}
\begin{itemize}[leftmargin=*,nosep]
\item \textbf{BMC} (9 blocks): Customer segments, value proposition, channels, customer relationships, revenue streams, key resources, key activities, key partnerships, cost structure.
\item \textbf{Product design}: Problem discovery $\to$ customer discovery $\to$ solution validation $\to$ build/launch.
\item \textbf{\hlpink{Lean validation}}: Don't spend until validated. Test. Learn. Iterate.
\item Problem--solution fit first. Go to market. Talk to customers. Don't assume.
\item BMC helps validate before building; fill canvas as you validate, not before.
\item \textbf{Market research}: Gartner, etc. for market forecasts, size, trends.
\item \textbf{Rebranding}: Expensive; may not fix underlying problem. Understand why it's not working first.
\end{itemize}

\section*{\color{eublue}Unit 5 -- Exploration vs.\ Exploitation}\handnote{Ambidextrous org.}
\begin{itemize}[leftmargin=*,nosep]
\item \textbf{\hlpink{Exploration}} (sensing) = Finding new opportunities. Searching.
\item \textbf{\hlpink{Exploitation}} (seizing) = Doing what works. Scaling. Capturing value.
\item Need both. Exploration finds; exploitation captures.
\item Microsoft: strong exploitation, weaker exploration.
\item \textbf{SDLC}: Plan $\to$ Build $\to$ Test $\to$ Deploy. Iterate.
\end{itemize}

\subsection*{Exploration vs.\ exploitation}
\begin{center}
\footnotesize
\begin{tabular}{@{}lll@{}}
\toprule
\textbf{Aspect} & \textbf{Exploration} & \textbf{Exploitation} \\
\midrule
Focus & Sensing, searching & Seizing, scaling \\
Activity & Find new & Do what works \\
\bottomrule
\end{tabular}
\end{center}

\section*{\color{eublue}Unit 6 -- Design Thinking}\handnote{Observe, don't ask.}
\begin{itemize}[leftmargin=*,nosep]
\item \hlpeach{Empathize $\to$ Define $\to$ Ideate $\to$ Prototype $\to$ Test}. Iterate.
\item Don't ask users what they want. Steve Jobs: ``People don't know what they want until you show them.''
\item \textbf{Observe behavior}. Ask about problems, past. Not solutions.
\item \textbf{\hlpink{Gemba}} = Go where work happens. Observe in real context.
\item \textbf{Two observers}: Reduces bias. Richer insight.
\item \textbf{\hlpink{Extreme users}}: Use product intensely. Reveal problems. IDEO cart: 22,000 kids/year injured.
\item Diverse teams (IDEO): engineers, psychologists, linguists. Different perspectives = better solutions.
\item Corporate trap: bosses reward desk time; best insight from the field.
\item \textit{\hlpeach{Exam tip}}: Experts in the process, not the product. Process transfers across domains.
\end{itemize}

\section*{\color{eublue}Unit 7 -- Innovation \& Technology}\handnote{S-curve. Disrupt.}
\begin{itemize}[leftmargin=*,nosep]
\item \textbf{\hlpink{S-curve}}: Slow start $\to$ rapid growth $\to$ plateau $\to$ decline.
\item \textbf{Investment pattern}: Early low; growth heavy; maturity shift to new tech.
\item \textbf{\hlpink{Disruptive innovation}}: New tech starts lower, cheaper. Improves. Overtakes old.
\item \textbf{Innovator's dilemma}: Focus on current customers. Miss new thing. Old tech still profitable.
\item \textbf{Kodak}: Dominant in film. Failed digital. Canon, Sony, Nikon won.\handnote{Classic case.}
\item \textbf{iPhone}: Apple didn't invent; revolutionized. Touchscreen, size, ease.
\item \textbf{First mover}: Advantage or disadvantage if too early.
\item Radical innovation: often from outside established companies.
\item \textbf{Uber}: Created new market. Disrupted taxi.
\item \textbf{Customer preferences}: May prefer new tech even if initially lower performance (convenience, speed).
\item \textbf{Examples}: Peloton, Tesla Cybertruck, Vision Pro, Threads, Clubhouse, BeReal.
\end{itemize}

\subsection*{\hlpink{S-curve phases}}
\begin{center}
\footnotesize
\begin{tabular}{@{}ll@{}}
\toprule
\textbf{Phase} & \textbf{Characteristic} \\
\midrule
Beginning & Slow, evolving, uncertain \\
Growth & Rapid improvement \\
Maturity & Plateau, incremental \\
Decline & Replaced by new tech \\
\bottomrule
\end{tabular}
\end{center}

\section*{\color{eublue}Unit 8 -- User Interviews \& Product Design}\handnote{Zoom or in-person. Behavioral questions.}
\begin{itemize}[leftmargin=*,nosep]
\item \textbf{\hlpink{User interviews}}: Main assignment. Interview real users (not friends in USA). Zoom or in-person; not text. Your notes fine. \textbf{Behavioral questions}: ``Tell me about the last time...''; past events, not hypotheticals. Avoid: ``What features would you want?'' Goal: how they describe problem, what they want, what frustrates them.
\item \textbf{Design thinking flow}: Empathize $\to$ Ideate $\to$ Prototype $\to$ Test. Ideation next; then prototype; final = something to show.
\item \textbf{\hlpink{Value engineering}}: Change design to cut cost or simplify. E.g.\ smaller spoon $\to$ less metal; sugar dispenser with holes $\to$ efficient portioning; plastic vs ceramic $\to$ cheaper, lighter.
\item \textbf{\hlpink{Platform design}}: One base, many variants. Shared components. E.g.\ BMW---same platform, many models.
\item \textbf{Mass customization}: Same platform, different configs for different segments. Economies of scale + variety.
\item \textbf{Less is more}: Remove features $\to$ cheaper, simpler. Users may not need every function. Simpler can win.
\end{itemize}\handnote{Value eng = cut cost. Platform = one base, many.}

\section*{\color{eublue}Quick revision -- Can you explain\ldots?}\handnote{Run through before exam.}
\begin{itemize}[leftmargin=*,nosep]
\item Value creation vs.\ capture. Porter value chain. Teece \& market access.
\item Takeda case (15K $\to$ 250M). Co-specialized assets.
\item Community before product. Willingness to pay. Don't build and hope.
\item Cost $\neq$ price. Fixed vs. variable. Economies of scale. Henry Ford.
\item BMC 9 blocks. Product design 4 stages. Lean validation.
\item Exploration vs.\ exploitation. Design thinking 5 stages.
\item Why not ask users? Gemba, extreme users, two observers.
\item S-curve, disruptive innovation, innovator's dilemma.
\item Kodak, iPhone, Uber, Peloton examples. First mover.
\item User interviews: behavioral questions, Zoom/in-person. Value engineering, platform design, mass customization.
\end{itemize}\handnote{If you can explain all these, you're ready.}

\section*{\hlpeach{Common question angles}}\handnote{Memorise these.}
\begin{itemize}[leftmargin=*,nosep]
\item \textbf{``Why create but not capture?''} Market access, distribution, patents, branding, scale. Takeda vs.\ Spanish research.
\item \textbf{``Apply design thinking.''} Empathize (observe) $\to$ define $\to$ ideate $\to$ prototype $\to$ test.
\item \textbf{``Exploration vs.\ exploitation?''} Both needed. Exploration finds; exploitation captures.
\item \textbf{``Why did Kodak fail?''} Innovator's dilemma; invested in film; didn't adapt; competitors did.
\item \textbf{``Apply BMC to [case].''} Customer segments, value prop, channels, key resources. Validate first.
\item \textbf{``Why might a product fail despite good design?''} No market fit; no willingness to pay; wrong community; didn't validate.
\item \textbf{``Cost vs.\ price?''} Cost = what you spend. Price = what you choose. Low cost $\neq$ low price.
\item \textbf{``Value engineering?''} Change design to cut cost or simplify (size, material, features). Sometimes less is more.
\item \textbf{``Platform design?''} One base, many variants. Shared components. Mass customization (e.g.\ BMW).
\end{itemize}

\end{multicols}

\newpage
\vspace*{0.3em}
\begin{center}
\textbf{\color{eublue}Quick recap -- last page}\\[0.3em]
\small\textcolor{notegrey}{(Section summaries, not margin notes)}
\end{center}
\vspace{0.4em}
\noindent\color{notegrey}
\setlength{\parskip}{0.3em}
\begin{multicols}{2}

\textbf{Unit 0 -- Big Ideas}\\
Creation $\neq$ capture. Great product not enough. Community before product.

\textbf{Unit 1 -- Value}\\
Creation vs.\ capture. Market access. Takeda. Teece. Co-specialized assets.

\textbf{Unit 2 -- Community}\\
Needs assessment. Willingness to pay. Don't build and hope. Testing as investing.

\textbf{Unit 3 -- Cost}\\
Cost $\neq$ price. Fixed vs.\ variable. Economies of scale. Henry Ford. Mass production.

\textbf{Unit 4 -- BMC \& Design}\\
9 blocks. Validate before build. Lean. Customer discovery.

\textbf{Unit 5 -- Exploration/Exploitation}\\
Sensing vs.\ seizing. Ambidextrous. SDLC.

\textbf{Unit 6 -- Design Thinking}\\
Observe don't ask. Gemba. Extreme users. 5 stages.

\textbf{Unit 7 -- Innovation}\\
S-curve. Disruptive. Kodak. Innovator's dilemma. First mover.

\textbf{Unit 8 -- Interviews \& Design}\\
User interviews (Zoom/in-person). Value engineering. Platform design. Mass customization. Less is more.

\textbf{Common angles}\\
Create vs capture. Design thinking. Kodak. BMC. Cost vs price. Willingness to pay.

\end{multicols}
\vfill
\normalsize\normalfont\color{black}
\end{document}
